% Original pages: 41-48
% Figures: none; one commutative diagram typeset directly
% Uncertainty log: none

\noindent\textbf{5.} Let $A[x]$ be the ring of polynomials in one indeterminate over a ring $A$.  Prove
that $A[x]$ is a flat $A$-algebra.  [Use Exercise 4.]

\medskip
\noindent\textbf{6.} For any $A$-module, let $M[x]$ denote the set of all polynomials in $x$ with coefficients in $M$, that is to say expressions of the form
\[
m_0+m_1x+\cdots+m_rx^r\qquad(m_i\in M).
\]
Defining the product of an element of $A[x]$ and an element of $M[x]$ in the obvious
way, show that $M[x]$ is an $A[x]$-module.

Show that $M[x]\cong A[x]\otimes_A M$.

\medskip
\noindent\textbf{7.} Let $\mathfrak p$ be a prime ideal in $A$.  Show that $\mathfrak p[x]$ is a prime ideal in $A[x]$.  If $\mathfrak m$ is a
maximal ideal in $A$, is $\mathfrak m[x]$ a maximal ideal in $A[x]$?

\medskip
\noindent\amstar\textbf{8.} i) If $M$ and $N$ are flat $A$-modules, then so is $M\otimes_A N$.

ii) If $B$ is a flat $A$-algebra and $N$ is a flat $B$-module, then $N$ is flat as an $A$-module.

\medskip
\noindent\amstar\textbf{9.} Let $0\longrightarrow M'\longrightarrow M\longrightarrow M''\longrightarrow0$ be an exact sequence of $A$-modules.  If $M'$ and
$M''$ are finitely generated, then so is $M$.

\medskip
\noindent\amstar\textbf{10.} Let $A$ be a ring, $\mathfrak a$ an ideal contained in the Jacobson radical of $A$; let $M$ be an
$A$-module and $N$ a finitely generated $A$-module, and let $u:M\longrightarrow N$ be a homomorphism.  If the induced homomorphism $M/\mathfrak aM\longrightarrow N/\mathfrak aN$ is surjective, then
$u$ is surjective.

\medskip
\noindent\textbf{11.} Let $A$ be a ring $\ne0$.  Show that $A^m\cong A^n\Rightarrow m=n$.

[Let $\mathfrak m$ be a maximal ideal of $A$ and let $\phi:A^m\longrightarrow A^n$ be an isomorphism.  Then
$1\otimes\phi:(A/\mathfrak m)\otimes A^m\longrightarrow(A/\mathfrak m)\otimes A^n$ is an isomorphism between vector spaces
of dimensions $m$ and $n$ over the field $k=A/\mathfrak m$.  Hence $m=n$.] (Cf. Chapter 3,
Exercise 15.)

If $\phi:A^m\longrightarrow A^n$ is surjective, then $m\geq n$.

If $\phi:A^m\longrightarrow A^n$ is injective, is it always the case that $m\leq n$?

\medskip
\noindent\textbf{12.} Let $M$ be a finitely generated $A$-module and $\phi:M\longrightarrow A^n$ a surjective homomorphism.  Show that $\operatorname{Ker}(\phi)$ is finitely generated.

[Let $e_1,\ldots,e_n$ be a basis of $A^n$ and choose $u_i\in M$ such that $\phi(u_i)=e_i$
$(1\leq i\leq n)$.  Show that $M$ is the direct sum of $\operatorname{Ker}(\phi)$ and the submodule
generated by $u_1,\ldots,u_n$.]

\medskip
\noindent\textbf{13.} Let $f:A\longrightarrow B$ be a ring homomorphism, and let $N$ be a $B$-module.  Regarding $N$
as an $A$-module by restriction of scalars, form the $B$-module $N_B=B\otimes_A N$.
Show that the homomorphism $g:N\longrightarrow N_B$ which maps $y$ to $1\otimes y$ is injective
and that $g(N)$ is a direct summand of $N_B$.

[Define $p:N_B\longrightarrow N$ by $p(b\otimes y)=by$, and show that $N_B=\operatorname{Im}(g)\oplus\operatorname{Ker}(p)$.]

\medskip
\noindent\emph{Direct limits}

\medskip
\noindent\amstar\textbf{14.} A partially ordered set $I$ is said to be a \emph{directed set} if for each pair $i,j$ in $I$ there
exists $k\in I$ such that $i\leq k$ and $j\leq k$.

Let $A$ be a ring, let $I$ be a directed set and let $(M_i)_{i\in I}$ be a family of $A$-modules
indexed by $I$.  For each pair $i,j$ in $I$ such that $i\leq j$, let $\mu_{ij}:M_i\longrightarrow M_j$ be an
$A$-homomorphism, and suppose that the following axioms are satisfied:
\begin{enumerate}
\renewcommand{\labelenumi}{(\arabic{enumi})}
\item $\mu_{ii}$ is the identity mapping of $M_i$, for all $i\in I$;
\item $\mu_{ik}=\mu_{jk}\circ\mu_{ij}$ whenever $i\leq j\leq k$.
\end{enumerate}
Then the modules $M_i$ and homomorphisms $\mu_{ij}$ are said to form a \emph{direct system}
$\mathbf M=(M_i,\mu_{ij})$ over the directed set $I$.

We shall construct an $A$-module $M$ called the \emph{direct limit} of the direct
system $\mathbf M$.  Let $C$ be the direct sum of the $M_i$, and identify each module $M_i$ with
its canonical image in $C$.  Let $D$ be the submodule of $C$ generated by all elements
of the form $x_i-\mu_{ij}(x_i)$ where $i\leq j$ and $x_i\in M_i$.  Let $M=C/D$, let $\mu:C\longrightarrow M$
be the projection and let $\mu_i$ be the restriction of $\mu$ to $M_i$.

The module $M$, or more correctly the pair consisting of $M$ and the family of
homomorphisms $\mu_i:M_i\longrightarrow M$, is called the \emph{direct limit} of the direct system $\mathbf M$,
and is written $\varinjlim M_i$.  From the construction it is clear that $\mu_i=\mu_j\circ\mu_{ij}$
whenever $i\leq j$.

\medskip
\noindent\amstar\textbf{15.} In the situation of Exercise 14, show that every element of $M$ can be written in
the form $\mu_i(x_i)$ for some $i\in I$ and some $x_i\in M_i$.

Show also that if $\mu_i(x_i)=0$ then there exists $j\geq i$ such that $\mu_{ij}(x_i)=0$
in $M_j$.

\medskip
\noindent\amstar\textbf{16.} Show that the direct limit is characterized (up to isomorphism) by the following
property.  Let $N$ be an $A$-module and for each $i\in I$ let $\alpha_i:M_i\longrightarrow N$ be an $A$-module homomorphism such that $\alpha_i=\alpha_j\circ\mu_{ij}$ whenever $i\leq j$.  Then there exists
a unique homomorphism $\alpha:M\longrightarrow N$ such that $\alpha_i=\alpha\circ\mu_i$ for all $i\in I$.

\medskip
\noindent\amstar\textbf{17.} Let $(M_i)_{i\in I}$ be a family of submodules of an $A$-module, such that for each pair of
indices $i,j$ in $I$ there exists $k\in I$ such that $M_i+M_j\subseteq M_k$.  Define $i\leq j$ to
mean $M_i\subseteq M_j$ and let $\mu_{ij}:M_i\longrightarrow M_j$ be the embedding of $M_i$ in $M_j$.  Show that
\[
\varinjlim M_i=\sum M_i=\bigcup M_i.
\]
In particular, any $A$-module is the direct limit of its finitely generated submodules.

\medskip
\noindent\amstar\textbf{18.} Let $\mathbf M=(M_i,\mu_{ij})$, $\mathbf N=(N_i,\nu_{ij})$ be direct systems of $A$-modules over the same
directed set.  Let $M,N$ be the direct limits and $\mu_i:M_i\longrightarrow M$, $\nu_i:N_i\longrightarrow N$ the
associated homomorphisms.

A \emph{homomorphism} $\Phi:\mathbf M\longrightarrow\mathbf N$ is by definition a family of $A$-module homomorphisms $\phi_i:M_i\longrightarrow N_i$ such that $\phi_j\circ\mu_{ij}=\nu_{ij}\circ\phi_i$ whenever $i\leq j$.  Show that
$\Phi$ defines a unique homomorphism $\phi=\varinjlim\phi_i:M\longrightarrow N$ such that $\phi\circ\mu_i=
\nu_i\circ\phi_i$ for all $i\in I$.

\medskip
\noindent\amstar\textbf{19.} A sequence of direct systems and homomorphisms
\[
\mathbf M\longrightarrow\mathbf N\longrightarrow\mathbf P
\]
is \emph{exact} if the corresponding sequence of modules and module homomorphisms
is exact for each $i\in I$.  Show that the sequence $M\longrightarrow N\longrightarrow P$ of direct limits is
then exact.  [Use Exercise 15.]

\medskip
\noindent\emph{Tensor products commute with direct limits}

\medskip
\noindent\amstar\textbf{20.} Keeping the same notation as in Exercise 14, let $N$ be any $A$-module.  Then
$(M_i\otimes N,\mu_{ij}\otimes1)$ is a direct system; let $P=\varinjlim(M_i\otimes N)$ be its direct limit.

For each $i\in I$ we have a homomorphism $\mu_i\otimes1:M_i\otimes N\longrightarrow M\otimes N$, hence by
Exercise 16 a homomorphism $\psi:P\longrightarrow M\otimes N$.  Show that $\psi$ is an isomorphism,
so that
\[
\varinjlim(M_i\otimes N)\cong(\varinjlim M_i)\otimes N.
\]
[For each $i\in I$, let $g_i:M_i\times N\longrightarrow M_i\otimes N$ be the canonical bilinear mapping.
Passing to the limit we obtain a mapping $g:M\times N\longrightarrow P$.  Show that $g$ is
$A$-bilinear and hence define a homomorphism $\phi:M\otimes N\longrightarrow P$.  Verify that
$\phi\circ\psi$ and $\psi\circ\phi$ are identity mappings.]

\medskip
\noindent\amstar\textbf{21.} Let $(A_i)_{i\in I}$ be a family of rings indexed by a directed set $I$, and for each pair $i\leq j$
in $I$ let $\alpha_{ij}:A_i\longrightarrow A_j$ be a ring homomorphism, satisfying conditions (1) and (2)
of Exercise 14.  Regarding each $A_i$ as a $\mathbb Z$-module we can then form the direct
limit $A=\varinjlim A_i$.  Show that $A$ inherits a ring structure from the $A_i$ so that the
mappings $A_i\longrightarrow A$ are ring homomorphisms.  The ring $A$ is the \emph{direct limit} of the
system $(A_i,\alpha_{ij})$.

If $A=0$ prove that $A_i=0$ for some $i\in I$.  [Remember that all rings have
identity elements!]

\medskip
\noindent\textbf{22.} Let $(A_i,\alpha_{ij})$ be a direct system of rings and let $\mathfrak N_i$ be the nilradical of $A_i$.  Show
that $\varinjlim\mathfrak N_i$ is the nilradical of $\varinjlim A_i$.

If each $A_i$ is an integral domain, then $\varinjlim A_i$ is an integral domain.

\medskip
\noindent\textbf{23.} Let $(B_\lambda)_{\lambda\in\Lambda}$ be a family of $A$-algebras.  For each finite subset $J$ of $\Lambda$ let $B_J$ denote
the tensor product (over $A$) of the $B_\lambda$ for $\lambda\in J$.  If $J'$ is another finite subset of $\Lambda$
and $J\subseteq J'$, there is a canonical $A$-algebra homomorphism $B_J\longrightarrow B_{J'}$.  Let $B$
denote the direct limit of the rings $B_J$ as $J$ runs through all finite subsets of $\Lambda$.
The ring $B$ has a natural $A$-algebra structure for which the homomorphisms
$B_J\longrightarrow B$ are $A$-algebra homomorphisms.  The $A$-algebra $B$ is the \emph{tensor product}
of the family $(B_\lambda)_{\lambda\in\Lambda}$.

\medskip
\noindent\emph{Flatness and Tor}

In these Exercises it will be assumed that the reader is familiar with the definition
and basic properties of the Tor functor.

\medskip
\noindent\amstar\textbf{24.} If $M$ is an $A$-module, the following are equivalent:

i) $M$ is flat;

ii) $\operatorname{Tor}_n^A(M,N)=0$ for all $n>0$ and all $A$-modules $N$;

iii) $\operatorname{Tor}_1^A(M,N)=0$ for all $A$-modules $N$.

[To show that (i) $\Rightarrow$ (ii), take a free resolution of $N$ and tensor it with $M$.  Since
$M$ is flat, the resulting sequence is exact and therefore its homology groups,
which are the $\operatorname{Tor}_n^A(M,N)$, are zero for $n>0$.  To show that (iii) $\Rightarrow$ (i), let
$0\longrightarrow N'\longrightarrow N\longrightarrow N''\longrightarrow0$ be an exact sequence.  Then, from the Tor exact
sequence,
\[
\operatorname{Tor}_1(M,N')\longrightarrow M\otimes N'\longrightarrow M\otimes N\longrightarrow M\otimes N''\longrightarrow0
\]
is exact.  Since $\operatorname{Tor}_1(M,N')=0$ it follows that $M$ is flat.]

\medskip
\noindent\amstar\textbf{25.} Let $0\longrightarrow N'\longrightarrow N\longrightarrow N''\longrightarrow0$ be an exact sequence, with $N''$ flat.  Then $N'$ is
flat $\Longleftrightarrow N$ is flat.  [Use Exercise 24 and the Tor exact sequence.]

\medskip
\noindent\amstar\textbf{26.} Let $N$ be an $A$-module.  Then $N$ is flat $\Longleftrightarrow\operatorname{Tor}_1(A/\mathfrak a,N)=0$ for all finitely
generated ideals $\mathfrak a$ in $A$.

[Show first that $N$ is flat if $\operatorname{Tor}_1(M,N)=0$ for all \emph{finitely generated} $A$-modules
$M$, by using \amref{2.19}.  If $M$ is finitely generated, let $x_1,\ldots,x_n$ be a set of generators
of $M$, and let $M_i$ be the submodule generated by $x_1,\ldots,x_i$.  By considering
the successive quotients $M_i/M_{i-1}$ and using Exercise 25, deduce that $N$ is flat
if $\operatorname{Tor}_1(M,N)=0$ for all \emph{cyclic} $A$-modules $M$, i.e., all $M$ generated by a single
element, and therefore of the form $A/\mathfrak a$ for some ideal $\mathfrak a$.  Finally use \amref{2.19} again
to reduce to the case where $\mathfrak a$ is a finitely generated ideal.]

\medskip
\noindent\textbf{27.} A ring $A$ is \emph{absolutely flat} if every $A$-module is flat.  Prove that the following are
equivalent:

i) $A$ is absolutely flat.

ii) Every principal ideal is idempotent.

iii) Every finitely generated ideal is a direct summand of $A$.

[i) $\Rightarrow$ ii).  Let $x\in A$.  Then $A/(x)$ is a flat $A$-module, hence in the diagram
\[
\begin{array}{ccc}
(x)\otimes A&\xrightarrow{\beta}&(x)\otimes A/(x)\\
\big\downarrow&&\big\downarrow\!\alpha\\
A&\longrightarrow&A/(x)
\end{array}
\]
the mapping $\alpha$ is injective.  Hence $\operatorname{Im}(\beta)=0$, hence $(x)=(x^2)$.  ii) $\Rightarrow$ iii).  Let
$x\in A$.  Then $x=ax^2$ for some $a\in A$, hence $e=ax$ is idempotent and we have
$(e)=(x)$.  Now if $e,f$ are idempotents, then $(e,f)=(e+f-ef)$.  Hence every
finitely generated ideal is principal, and generated by an idempotent $e$, hence is a
direct summand because $A=(e)\oplus(1-e)$.  iii) $\Rightarrow$ i).  Use the criterion of
Exercise 26.]

\medskip
\noindent\textbf{28.} A Boolean ring is absolutely flat.  The ring of Chapter 1, Exercise 7 is absolutely
flat.  Every homomorphic image of an absolutely flat ring is absolutely flat.  If a
local ring is absolutely flat, then it is a field.

If $A$ is absolutely flat, every non-unit in $A$ is a zero-divisor.

\section{Rings and Modules of Fractions}

The formation of rings of fractions and the associated process of localization
are perhaps the most important technical tools in commutative algebra.  They
correspond in the algebro-geometric picture to concentrating attention on an
open set or near a point, and the importance of these notions should be self-evident.  This chapter gives the definitions and simple properties of the formation
of fractions.

The procedure by which one constructs the rational field $\mathbb Q$ from the ring
of integers $\mathbb Z$ (and embeds $\mathbb Z$ in $\mathbb Q$) extends easily to any integral domain $A$ and
produces the \emph{field of fractions} of $A$.  The construction consists in taking all
ordered pairs $(a,s)$ where $a,s\in A$ and $s\ne0$, and setting up an equivalence
relation between such pairs:
\[
(a,s)\equiv(b,t)\Longleftrightarrow at-bs=0.
\]
This works only if $A$ is an integral domain, because the verification that the
relation is transitive involves canceling, i.e. the fact that $A$ has no zero-divisor
$\ne0$.  However, it can be generalized as follows:

Let $A$ be any ring.  A \emph{multiplicatively closed subset} of $A$ is a subset $S$ of $A$
such that $1\in S$ and $S$ is closed under multiplication: in other words $S$ is a subsemigroup of the multiplicative semigroup of $A$.  Define a relation $\equiv$ on $A\times S$
as follows:
\[
(a,s)\equiv(b,t)\Longleftrightarrow(at-bs)u=0\text{ for some }u\in S.
\]
Clearly this relation is reflexive and symmetric.  To show that it is transitive,
suppose $(a,s)\equiv(b,t)$ and $(b,t)\equiv(c,u)$.  Then there exist $v,w$ in $S$ such that
$(at-bs)v=0$ and $(bu-ct)w=0$.  Eliminate $b$ from these two equations
and we have $(au-cs)tvw=0$.  Since $S$ is closed under multiplication, we have
$tvw\in S$, hence $(a,s)\equiv(c,u)$.  Thus we have an equivalence relation.  Let $a/s$
denote the equivalence class of $(a,s)$, and let $S^{-1}A$ denote the set of equivalence
classes.  We put a ring structure on $S^{-1}A$ by defining addition and multiplication
of these ``fractions'' $a/s$ in the same way as in elementary algebra: that is,
\[
(a/s)+(b/t)=(at+bs)/st,
\]
\[
(a/s)(b/t)=ab/st.
\]

\emph{Exercise.  Verify that these definitions are independent of the choices of representatives $(a,s)$ and $(b,t)$, and that $S^{-1}A$ satisfies the axioms of a commutative ring with identity.}

We also have a ring homomorphism $f:A\longrightarrow S^{-1}A$ defined by $f(x)=x/1$.
This is not in general injective.

\emph{Remark.}  If $A$ is an integral domain and $S=A-\{0\}$, then $S^{-1}A$ is the field
of fractions of $A$.

The ring $S^{-1}A$ is called the \emph{ring of fractions} of $A$ with respect to $S$.  It has a
\emph{universal property}:

\medskip
\noindent\textit{\textbf{Proposition 3.1.\amtarget{3.1}} Let $g:A\longrightarrow B$ be a ring homomorphism such that $g(s)$
is a unit in $B$ for all $s\in S$.  Then there exists a unique ring homomorphism
$h:S^{-1}A\longrightarrow B$ such that $g=h\circ f$.}

\noindent\emph{Proof.}  i) Uniqueness.  If $h$ satisfies the conditions, then $h(a/1)=hf(a)=g(a)$
for all $a\in A$; hence, if $s\in S$,
\[
h(1/s)=h((s/1)^{-1})=h(s/1)^{-1}=g(s)^{-1}
\]
and therefore $h(a/s)=h(a/1)\cdot h(1/s)=g(a)g(s)^{-1}$, so that $h$ is uniquely
determined by $g$.

ii) Existence.  Let $h(a/s)=g(a)g(s)^{-1}$.  Then $h$ will clearly be a ring homomorphism provided that it is well-defined.  Suppose then that $a/s=a'/s'$; then
there exists $t\in S$ such that $(as'-a's)t=0$, hence
\[
(g(a)g(s')-g(a')g(s))g(t)=0;
\]
now $g(t)$ is a unit in $B$, hence $g(a)g(s)^{-1}=g(a')g(s')^{-1}$.\proofbox

The ring $S^{-1}A$ and the homomorphism $f:A\longrightarrow S^{-1}A$ have the following
properties:

1) $s\in S\Rightarrow f(s)$ is a unit in $S^{-1}A$;

2) $f(a)=0\Rightarrow as=0$ for some $s\in S$;

3) Every element of $S^{-1}A$ is of the form $f(a)f(s)^{-1}$ for some $a\in A$ and some
$s\in S$.

Conversely, these three conditions determine the ring $S^{-1}A$ up to isomorphism.  Precisely:

\medskip
\noindent\textit{\textbf{Corollary 3.2.\amtarget{3.2}} If $g:A\longrightarrow B$ is a ring homomorphism such that}

\noindent\textit{i) $s\in S\Rightarrow g(s)$ is a unit in $B$;}

\noindent\textit{ii) $g(a)=0\Rightarrow as=0$ for some $s\in S$;}

\noindent\textit{iii) Every element of $B$ is of the form $g(a)g(s)^{-1}$; then there is a unique
isomorphism $h:S^{-1}A\longrightarrow B$ such that $g=h\circ f$.}

\noindent\emph{Proof.}  By \amref{3.1} we have to show that $h:S^{-1}A\longrightarrow B$, defined by
\[
h(a/s)=g(a)g(s)^{-1}
\]
(this definition uses i)) is an isomorphism.  By iii), $h$ is surjective.  To show $h$ is
injective, look at the kernel of $h$: if $h(a/s)=0$, then $g(a)=0$, hence by ii) we
have $at=0$ for some $t\in S$, hence $(a,s)\equiv(0,1)$, i.e., $a/s=0$ in $S^{-1}A$.\proofbox

\medskip
\noindent\textbf{Examples.}  1) Let $\mathfrak p$ be a prime ideal of $A$.  Then $S=A-\mathfrak p$ is multiplicatively
closed (in fact $A-\mathfrak p$ is multiplicatively closed $\Longleftrightarrow\mathfrak p$ is prime).  We write $A_\mathfrak p$ for
$S^{-1}A$ in this case.  The elements $a/s$ with $a\in\mathfrak p$ form an ideal $\mathfrak m$ in $A_\mathfrak p$.  If $b/t\notin\mathfrak m$,
then $b\notin\mathfrak p$, hence $b\in S$ and therefore $b/t$ is a unit in $A_\mathfrak p$.  It follows that if $\mathfrak a$ is an
ideal in $A_\mathfrak p$ and $\mathfrak a\nsubseteq\mathfrak m$, then $\mathfrak a$ contains a unit and is therefore the whole ring.
Hence $\mathfrak m$ is the only maximal ideal in $A_\mathfrak p$; in other words, $A_\mathfrak p$ is a \emph{local ring}.

The process of passing from $A$ to $A_\mathfrak p$ is called \emph{localization at $\mathfrak p$}.

2) $S^{-1}A$ is the zero ring $\Longleftrightarrow0\in S$.

3) Let $f\in A$ and let $S=\{f^n\}_{n\geq0}$.  We write $A_f$ for $S^{-1}A$ in this case.

4) Let $\mathfrak a$ be any ideal in $A$, and let $S=1+\mathfrak a={}$ set of all $1+x$ where
$x\in\mathfrak a$.  Clearly $S$ is multiplicatively closed.

5) Special cases of 1) and 3):

i) $A=\mathbb Z$, $\mathfrak p=(p)$, $p$ a prime number; $A_\mathfrak p={}$ set of all rational numbers
$m/n$ where $n$ is prime to $p$; if $f\in\mathbb Z$ and $f\ne0$, then $A_f$ is the set of all rational
numbers whose denominator is a power of $f$.

ii) $A=k[t_1,\ldots,t_n]$, where $k$ is a field and the $t_i$ are independent indeterminates, $\mathfrak p$ a prime ideal in $A$.  Then $A_\mathfrak p$ is the ring of all rational functions $f/g$,
where $g\notin\mathfrak p$.  If $V$ is the variety defined by the ideal $\mathfrak p$, that is to say the set of all
$x=(x_1,\ldots,x_n)\in k^n$ such that $f(x)=0$ whenever $f\in\mathfrak p$, then (provided $k$ is
infinite) $A_\mathfrak p$ can be identified with the ring of all rational functions on $k^n$ which
are defined at almost all points of $V$; it is the local ring of $k^n$ \emph{along the variety $V$}.
This is the prototype of the local rings which arise in algebraic geometry.

The construction of $S^{-1}A$ can be carried through with an $A$-module $M$ in
place of the ring $A$.  Define a relation $\equiv$ on $M\times S$ as follows:
\[
(m,s)\equiv(m',s')\Longleftrightarrow\exists t\in S\text{ such that }t(sm'-s'm)=0.
\]
As before, this is an equivalence relation.  Let $m/s$ denote the equivalence class
of the pair $(m,s)$, let $S^{-1}M$ denote the set of such fractions, and make $S^{-1}M$
into an $S^{-1}A$-module with the obvious definitions of addition and scalar
multiplication.  As in Examples 1) and 3) above, we write $M_\mathfrak p$ instead of $S^{-1}M$
when $S=A-\mathfrak p$ ($\mathfrak p$ prime) and $M_f$ when $S=\{f^n\}_{n\geq0}$.

Let $u:M\longrightarrow N$ be an $A$-module homomorphism.  Then it gives rise to an
$S^{-1}A$-module homomorphism $S^{-1}u:S^{-1}M\longrightarrow S^{-1}N$, namely $S^{-1}u$ maps $m/s$
to $u(m)/s$.  We have $S^{-1}(v\circ u)=(S^{-1}v)\circ(S^{-1}u)$.

\medskip
\noindent\textit{\textbf{Proposition 3.3.\amtarget{3.3}} The operation $S^{-1}$ is exact, i.e., if $M'\xrightarrow{f}M\xrightarrow{g}M''$ is
exact at $M$, then $S^{-1}M'\xrightarrow{S^{-1}f}S^{-1}M\xrightarrow{S^{-1}g}S^{-1}M''$ is exact at $S^{-1}M$.}

\noindent\emph{Proof.}  We have $g\circ f=0$, hence $S^{-1}g\circ S^{-1}f=S^{-1}(0)=0$, hence $\operatorname{Im}(S^{-1}f)
\subseteq\operatorname{Ker}(S^{-1}g)$.  To prove the reverse inclusion, let $m/s\in\operatorname{Ker}(S^{-1}g)$, then
$g(m)/s=0$ in $S^{-1}M''$, hence there exists $t\in S$ such that $tg(m)=0$ in $M''$.  But
$tg(m)=g(tm)$ since $g$ is an $A$-module homomorphism, hence $tm\in\operatorname{Ker}(g)=
\operatorname{Im}(f)$ and therefore $tm=f(m')$ for some $m'\in M'$.  Hence in $S^{-1}M$ we have
$m/s=f(m')/st=(S^{-1}f)(m'/st)\in\operatorname{Im}(S^{-1}f)$.  Hence $\operatorname{Ker}(S^{-1}g)\subseteq\operatorname{Im}(S^{-1}f)$.\proofbox

In particular, it follows from \amref{3.3} that if $M'$ is a submodule of $M$, the mapping $S^{-1}M'\longrightarrow S^{-1}M$ is injective and therefore $S^{-1}M'$ can be regarded as a
submodule of $S^{-1}M$.  With this convention,

\medskip
\noindent\textit{\textbf{Corollary 3.4.\amtarget{3.4}} Formation of fractions commutes with formation of finite
sums, finite intersections and quotients.  Precisely, if $N,P$ are submodules
of an $A$-module $M$, then}

\noindent\textit{i) $S^{-1}(N+P)=S^{-1}(N)+S^{-1}(P)$}

\noindent\textit{ii) $S^{-1}(N\cap P)=S^{-1}(N)\cap S^{-1}(P)$}

\noindent\textit{iii) the $S^{-1}A$-modules $S^{-1}(M/N)$ and $(S^{-1}M)/(S^{-1}N)$ are isomorphic.}

\noindent\emph{Proof.}  i) follows readily from the definitions and ii) is easy to veri{}fy:
if $y/s=z/t$ ($y\in N$, $z\in P$, $s,t\in S$) then $u(ty-sz)=0$ for some $u\in S$, hence
$w=uty=usz\in N\cap P$ and therefore $y/s=w/stu\in S^{-1}(N\cap P)$.  Consequently
$S^{-1}N\cap S^{-1}P\subseteq S^{-1}(N\cap P)$, and the reverse inclusion is obvious.

iii) Apply $S^{-1}$ to the exact sequence $0\longrightarrow N\longrightarrow M\longrightarrow M/N\longrightarrow0$.\proofbox

\medskip
\noindent\textit{\textbf{Proposition 3.5.\amtarget{3.5}} Let $M$ be an $A$-module.  Then the $S^{-1}A$ modules $S^{-1}M$ and
$S^{-1}A\otimes_A M$ are isomorphic; more precisely, there exists a unique isomorphism $f:S^{-1}A\otimes_A M\longrightarrow S^{-1}M$ for which}
\begin{equation}
f((a/s)\otimes m)=am/s\text{ for all }a\in A,\ m\in M,\ s\in S.
\tag{1}
\end{equation}
\noindent\emph{Proof.}  The mapping $S^{-1}A\times M\longrightarrow S^{-1}M$ defined by
\[
(a/s,m)\longmapsto am/s
\]
is $A$-bilinear, and therefore by the universal property \amref{2.12} of the tensor product
induces an $A$-homomorphism
\[
f:S^{-1}A\otimes_A M\longrightarrow S^{-1}M
\]
satisfying (1).  Clearly $f$ is surjective, and is uniquely defined by (1).

Let $\sum_i(a_i/s_i)\otimes m_i$ be any element of $S^{-1}A\otimes M$.  If $s=\prod_i s_i\in S$,
$t_i=\prod_{j\ne i}s_j$, we have
\[
\sum_i\frac{a_i}{s_i}\otimes m_i
=\sum_i\frac{a_it_i}{s}\otimes m_i
=\sum_i\frac1s\otimes a_it_im_i
=\frac1s\otimes\sum_i a_it_im_i,
\]
